\section{Open questions}\label{sec:open}

Two questions are left open by this paper, and this section states them
precisely enough that answering either is a definite piece of work. A third
question, on the closed form of a constant arising in the Borel plane, belonged
to an earlier scope for this paper and has left it together with the material
it depended on; it is not restated here, because stating it would require
machinery that no section above sets up.

\subsection{Uniqueness of the normalised representative}\label{ssec:gauge}

The membership argument of \S\ref{ssec:verdict} compares System
$\mathsf C$ with a member of the van der Put--Saito family. That comparison is
between \emph{representatives}, and the family is presented in a normalised
form. It is natural to ask whether the normalisation determines the
representative --- whether, inside a projective class, van der Put--Saito's
normalised form is reached by exactly one module. If it were, membership would
follow from a projective match alone, and the whole of \S\ref{ssec:verdict}'s
second step would be unnecessary.

The obvious argument for uniqueness is that a trace-free representative is
unique. That argument is false, and the counterexample is explicit: the
$2$-torsion object $L=\partial-\tfrac{1}{2z}$ of \S\ref{ssec:tracefree}
(\texttt{S17-02}, \textbf{PROVEN}) is trivial after tensoring with itself but
not before, so tensoring by it preserves both the projective class and the
triviality of the determinant while changing the module. A projective class can
therefore contain more than one trace-free module up to isomorphism, and
``exactly one trace-free representative per projective class'' is false as a
general statement.

What replaces it is not settled here.

\begin{remark}[The question, deliberately not answered; \texttt{S17-07}, \textbf{CONJECTURED}]
Since trace-freeness alone is insufficient, the correct sharper hypothesis
under which a projective class has a unique representative in van der
Put--Saito's normalised form is not settled. A $2$-torsion twist shifts local
exponents by $\tfrac12$, but individual local exponents depend on the choice of
lattice, so whether the normalisation ``local exponents $\pm\theta_i/2$''
excludes such twists needs an argument this work does not make. It is reported
as an open question rather than guessed.
\end{remark}

Two remarks on the status of this gap.

It is not load-bearing. The membership result does not depend on uniqueness,
precisely because the comparison of $(P,Q)$ was carried out directly rather
than deduced from a projective match. The question is open, and the paper is
not waiting on it.

It is also not closed by the Galois symmetry that appears elsewhere in this
work, and the point is worth making explicitly, because two statements in this
paper carry a $\sqrt3$ and a reader could otherwise join them.

\begin{remark}[The elimination, and what its symmetry is about; \texttt{S25-02}, \textbf{PROVEN}]
The ideal generated by the six matching generators, saturated at $q(q-1)\neq0$,
is zero-dimensional of degree exactly $2$ over $K=\Q(i\sqrt{11})$. Its lex
eliminant in $\theta_\infty$ is proportional to
$27\theta_\infty^{2}+108\theta_\infty+107$ --- exactly the polynomial
registered as a prediction before the computation was run --- with roots
$-2-\sqrt3/9$ and $-2+\sqrt3/9$. The two solutions are Galois conjugate over
$K$ under $\sqrt3\mapsto-\sqrt3$.
\end{remark}

That conjugacy relates the \emph{two solutions of the elimination}, and hence
the two branches of the sign pair, both of which were run
(\S\ref{ssec:numerical-scope}). It is a statement about parameters. The open
question of this subsection is about a $z$-dependent gauge transformation
between modules, whose effect is not a substitution in a parameter field. The
two are unrelated, and neither bears on the other: \textbf{the
$\sqrt3\mapsto-\sqrt3$ conjugacy does not close the gauge question.} It is
stated here so that the reader does not have to decide for themselves whether a
connection was intended.

\subsection{Independence of the monodromy convention}\label{ssec:convention-independence}

The comparison of characters in \S\ref{sec:monodromy} is carried out in one
convention, selected from sixteen by four pre-registered conditions
(\S\ref{ssec:numerical-scope}, \texttt{S22-11}). The selection is a result: the
conditions were fixed first, and fifteen candidates fail them quantitatively.
What is open is the next statement along, which is stronger and which this
paper does not make.

\begin{quote}
\emph{Question.} Is the identification of System $\mathsf C$'s point in the
character variety independent of the four conditions used to fix the
convention --- that is, would a different, equally defensible set of
conventions yield a point equivalent to this one under the corresponding
change of coordinates?
\end{quote}

The distinction matters because the two statements are easy to conflate. What
is established is that, \emph{given} the four conditions, the convention is
forced and the comparison succeeds. What is not established is that the
conditions themselves are inessential. A convention is a choice of orientation,
of side for the group action, and of which loop is which; a different choice
induces a computable change of coordinates on the character variety, and
showing independence means exhibiting that change and verifying that it carries
the located point to the located point.

Answering it is mechanical rather than deep: enumerate the conventions as
before, and for each, transport the located point along the induced coordinate
change instead of discarding it. The reason it was not done is that fifteen of
the sixteen do not satisfy the conditions, so the machinery built for this work
excludes them before any transport is attempted. A study of independence would
need the excluded fifteen carried forward rather than dropped, which is a
different computation from the one performed.

Until it is done, the membership statement should be read with the scope given
in \S\ref{ssec:numerical-scope}: confirmed for the trace coordinates
$x_1,x_2,x_3$, under the one convention of sixteen that passed.
